\section{Introduction}
\label{sec:intro}

Should language models think longer to answer harder questions? The prevailing wisdom in chain-of-thought (\CoT) reasoning research suggests yes: longer, more detailed reasoning traces tend to improve accuracy on challenging benchmarks~\citep{wei2022chainofthought,wang2022selfconsistency}. This assumption underpins the design of recent reasoning-focused models and inference-time scaling strategies~\citep{aggarwal2025l1}. Yet a growing body of evidence suggests the relationship between reasoning depth and performance is more nuanced---particularly when models encounter inputs that differ from their training distribution.

{\bf Why does trace length matter for robustness?} As LLMs are deployed in high-stakes reasoning tasks spanning medical diagnosis, legal analysis, and scientific discovery, their reliability under distribution shift becomes critical. Practitioners must decide how much inference-time compute to allocate: too little and the model may fail on complex problems; too much and the model may ``overthink,'' accumulating errors and hallucinations that degrade performance~\citep{su2025underthinkingoverthinking,shojaee2025illusion}. Recent work has shown that longer \CoT traces can hurt accuracy on both easy~\citep{wu2025whenmoreisless} and out-of-distribution tasks~\citep{li2025mirage}, but no study has systematically isolated the causal effect of trace length on OOD robustness under controlled conditions.

{\bf What is missing?} Prior work has studied the length--accuracy tradeoff~\citep{wu2025whenmoreisless,jin2024impact} and OOD fragility of \CoT~\citep{li2025mirage,wang2024nora} as separate phenomena. Length-controlled studies typically evaluate only in-distribution performance~\citep{aggarwal2025l1,lee2025tokencomplexity}, while robustness studies rarely control for trace length. The interaction between these two dimensions---how trace length modulates robustness---remains an open question. Furthermore, adaptive trace-length policies conditioned on model uncertainty have been proposed conceptually~\citep{zhang2025reasoningbudget} but not compared against strong fixed-length baselines under matched evaluation conditions.

{\bf Our approach.} We design a controlled experiment that isolates reasoning trace length as the sole independent variable. We define five prompt-based policies---\pnone (direct answer), \pshort (1--2 steps), \pmedium (4--6 steps), \plong (8+ steps), and \padaptive (short with confidence-triggered retry)---and evaluate each on 250 questions spanning five benchmarks from in-distribution arithmetic (\gsmk) to far-OOD commonsense reasoning (\csqa). All experiments use deterministic decoding (temperature $= 0$) on \nano, with a confirmatory replication on \mini.

Our key findings are striking. In-distribution accuracy follows an inverted U-shape, peaking at \pshort/\pmedium (90\%) and declining at \plong (80\%). OOD performance degrades monotonically with trace length: \pshort achieves 46\% on \mathfive while \plong achieves only 40\%, a gap that widens to 70\% vs.\ 15\% on the stronger \mini model. The \padaptive policy achieves the best OOD-near accuracy (56\%) and the smallest robustness gap (30 percentage points), while a $70\times$ efficiency advantage of concise over verbose reasoning highlights the cost of overthinking.

We make the following contributions:
\begin{itemize}[leftmargin=*,itemsep=0pt,topsep=0pt]
    \item We conduct the first controlled study of reasoning trace length as a causal factor in OOD robustness, demonstrating a strong non-monotonic relationship across five benchmarks and two model sizes.
    \item We show that verbose reasoning (\plong) catastrophically degrades OOD performance---by up to 55 percentage points on \mini---providing empirical support for the ``overthinking'' hypothesis.
    \item We evaluate an uncertainty-adaptive trace-length policy that achieves the best OOD-near accuracy and smallest robustness gap, while identifying confidence calibration as its key bottleneck.
    \item We quantify the cost--accuracy--robustness tradeoff, showing that short reasoning dominates long reasoning in token-normalized efficiency by up to $70\times$.
\end{itemize}

\para{Paper organization.} \Secref{sec:related} positions this work against prior research. \Secref{sec:method} details the experimental protocol. \Secref{sec:results} reports quantitative findings. \Secref{sec:discussion} discusses implications and limitations. \Secref{sec:conclusion} concludes.
